\section{结论边界：General Fallback，而不是 Universal Winner}

前面跨越图像、bytes、multimodal autoencoding 和 optical flow，本节回到正确的适用范围。讲者把 Perceiver 定位为：当我们不知道正确结构、希望快速接入新数据形态或需要共享多模态接口时，一个很强的 general fallback。若结构已知、数据少或 latency 极严，专用偏置仍可能更优。

\lecturefigure{slide-37-general-perception-summary.jpg}{讲座总结：线性输入 scaling、drop-in Transformer interface、多模态与结构化输出。}{Stanford Online 当前官方视频 00:56:14--00:56:21。}

\begin{knowledgebox}{读图：把总结改写成条件句}
Perceiver 对 input size 线性，是在 latent count 固定时；可替换 Transformer，是指 attention-based interface，而非所有训练 recipe 不变；多模态统一，是架构层面可共享，而非本讲已完成所有任务联合训练。
\end{knowledgebox}

\teachervoice{结尾 Q\&A 给出最重要的负面边界：团队还不擅长 small-scale data，因为弱偏置需要从数据中恢复结构；真正把多种模态和任务联合训练也仍是未来工作。}

相关工作页把方法放回历史：multimodal learning、efficient attention、tokenization-free language、optical flow、Transformer 与 iterative attention 都贡献了组件。阅读这些来源可以避免把 Perceiver 当作凭空出现的单一技巧。

\lecturefigure{slide-38-references.jpg}{讲座相关工作：multimodal、efficient attention、language、optical flow 与 iterative attention。}{Stanford Online 当前官方视频 00:56:22--00:58:49。}

\begin{importantbox}{选择架构的决策顺序}
先问输入/输出结构是否已知，再问数据量是否足以学习结构，然后比较 domain-specific baseline、Perceiver 的 latent/output query 配置、真实硬件速度与跨域复用成本。不要从“统一”二字直接跳到替换全部专用模型。
\end{importantbox}

最后一页给出 Perceiver 和 Perceiver IO 两篇核心论文及共同作者。它不仅是致谢页，也是复现入口：架构图说明接口，原论文补齐训练细节、数据、ablation 和任务实现。

\lecturefigure{slide-39-thank-you-and-papers.jpg}{两篇 Perceiver 核心论文、共同作者与讲座结束页。}{Stanford Online 当前官方视频 00:58:50--00:58:55。}

\begin{warningbox}{两个最常见的过度结论}
第一，“linear scaling”不等于端到端永远线性或更快；第二，“any data as a table”不等于任务结构和数据治理消失。Latent 数、output 数、训练数据、位置编码和硬件 kernel 仍决定可用性。
\end{warningbox}

\subsection{本章小结}

Perceiver 是结构不明确、新模态和统一接口场景的强起点；在已知领域、小数据和严格延迟场景，专用偏置仍占优势。它的研究价值在于把 generality-speed tradeoff 变成可调的 latent/query interface。

